%%%%%%%%%%%%%%%%%%%%%%%%%%%%%%%%%%%%%%%%%
% Plasmati Graduate CV
% LaTeX Template
% Version 1.0 (24/3/13)
%
% This template has been downloaded from:
% http://www.LaTeXTemplates.com
%
% Original author:
% Alessandro Plasmati (alessandro.plasmati@gmail.com)
%
% License:
% CC BY-NC-SA 3.0 (http://creativecommons.org/licenses/by-nc-sa/3.0/)
%
% Important note:
% This template needs to be compiled with XeLaTeX.
% The main document font is called Fontin and can be downloaded for free
% from here: http://www.exljbris.com/fontin.html
%
%%%%%%%%%%%%%%%%%%%%%%%%%%%%%%%%%%%%%%%%%

%----------------------------------------------------------------------------------------
%	PACKAGES AND OTHER DOCUMENT CONFIGURATIONS
%----------------------------------------------------------------------------------------

\documentclass[a4paper,11pt]{article} % Default font size and paper size

\usepackage{tabularx} % For enchanced Tabulars
\usepackage{fontspec} % For loading fonts
\defaultfontfeatures{Mapping=tex-text}
\setmainfont[SmallCapsFont = Fontin SmallCaps]{Fontin} % Main document font

\usepackage{xunicode,xltxtra,url,parskip} % Formatting packages

\usepackage[usenames,dvipsnames]{xcolor} % Required for specifying custom colors

\usepackage{tikz}
\usetikzlibrary{calc}
\usepackage{multirow}
\usepackage{wrapfig}
\usepackage{fontawesome}

\usepackage[big]{layaureo} % Margin formatting of the A4 page, an alternative to layaureo can be \usepackage{fullpage}
% To reduce the height of the top margin uncomment: 
\addtolength{\oddsidemargin}{-0.65in}
\addtolength{\evensidemargin}{-0.65in}
\addtolength{\textwidth}{1.3in}

\addtolength{\topmargin}{-0.3in}
\addtolength{\textheight}{1in}

\usepackage{hyperref} % Required for adding links	and customizing them
\definecolor{linkcolour}{rgb}{0,0.2,0.6} % Link color
\hypersetup{colorlinks,breaklinks,urlcolor=linkcolour,linkcolor=linkcolour} % Set link colors throughout the document

\usepackage{titlesec} % Used to customize the \section command
\titleformat{\section}{\Large\scshape\raggedright}{}{0em}{}[\titlerule] % Text formatting of sections
\titlespacing{\section}{0pt}{3pt}{3pt} % Spacing around sections

\titleformat{\subsection}{\normalsize\scshape\raggedright}{}{0em}{}[\titlerule] % Text formatting of sections
\titlespacing{\subsection}{0pt}{3pt}{3pt} % Spacing around sections

\begin{document}

\pagestyle{empty} % Removes page numbering

\font\fb=''[cmr10]'' % Change the font of the \LaTeX command under the skills section

%\begin{tikzpicture}[remember picture, overlay]
 % \draw[line width = 1.5pt, style = double] (-0.25in,0.45in) rectangle
%  ($(current page.south east) + (-0.25in,0.25in)$);
%\end{tikzpicture}

%----------------------------------------------------------------------------------------
%	NAME AND CONTACT INFORMATION
%----------------------------------------------------------------------------------------

\par{
\centering{
\Huge Tristan \textsc{Overney} \\
\Large \textit{ }
}
\bigskip
\par
} % Your name

\begin{wrapfigure}{r}{2.1in}
\vspace{-3.5cm}
\footnotesize

\begin{flushright}

\hspace{1cm}
\href{mailto:tristan@overney.org}{tristan@overney.org} 
\hspace{0.1cm}
\faEnvelope{}
\\ 

\href{https://www.linkedin.com/in/toverney/en}{ch.linkedin.com/in/toverney/en}
\hspace{0.1cm}
\faLinkedin{}
\\

\href{https://github.com/toverney}{github.com/toverney}
\hspace{0.1cm} 
\faGithub{}
\\


\end{flushright}

\end{wrapfigure}

%----------------------------------------------------------------------------------------
%	EDUCATION
%----------------------------------------------------------------------------------------

\section{Education}

\begin{tabularx}{\textwidth}{r X r} 	
2015--\textsc{Present} & 
\textbf{Masters Degree in Computer Science} & 
\textit{Swiss Federal Institute of Technology (EPFL)}\\

%------------------------------------------------

2011--2015 & 
\textbf{Bachelor Degree in Computer Science} & 
\textit{ EPFL }\\

%------------------------------------------------

\end{tabularx}

%----------------------------------------------------------------------------------------
%	WORK EXPERIENCE & PROJECTS 
%----------------------------------------------------------------------------------------

\section{Work Experience and Projects}

\begin{tabularx}{\textwidth}{r|X}

\textsc{July 2016-Sept 2016} & \textbf{Datamining for comparative analysis}\\
\multirow{4}{*}{ \includegraphics[width=50pt]{img/nexthinkLogo.png}}
& \emph{Nexthink R\&D}\\
& \footnotesize{With Nexthink's datas and using Spark (ml and sql), I implemented a module that runs a custom random forest and find rules to identify the machines that have issues (i.e. rules that explain the cause of the problem) and then modify their application to integrate my new interface. This allows the end-users to directly see what could be the factor causing issue to the service instead of having to investigate manually}\\
\multicolumn{2}{c}{} \\

%------------------------------------------------

\textsc{Feb 2015-Current} & \textbf{Teaching Assistant for Scala courses and MOOCs} \\
\multirow{4}{*}{ \includegraphics[width=50pt]{img/Scala.eps}}
& \emph{Courses taught by Martin Odersky and Viktor Kuncak at EPFL}\\ 
& \footnotesize{Through different courses (Functional, Reactive, Parallel or Concurrent programming) I am helping students discovering the Scala programming language and its paradigms by answering their questions. I also wrote and corrected exam questions, helped with the grading infrastructure, the assignments and with the EPFL Scala MOOCs on Coursera (forum and setting up) }\\
\multicolumn{2}{c}{} \\

%------------------------------------------------

\textsc{Sept 2015-Jan 2016} & \textbf{Scala.js Compiler Optimisations} \\
\multirow{4}{*}{ \includegraphics[width=50pt]{img/scalajs.png}}
& \emph{Performance and Optimisation on Scala.js}\\ 
& \footnotesize{Working with a partner, we first worked on a
benchmarking framework for Scala.js project, the JavaScript implementation of
the Scala compiler, by porting the existing Scalameter DSL.
We then implemented compiler optimisations such as constant folding and
constructor inlining.}\\
\multicolumn{2}{c}{} \\

%------------------------------------------------

\textsc{Sept-Dec 2014} & \textbf{Software Engineering Course Project} \\
\multirow{4}{*}{ \includegraphics[width=50pt]{img/epfl.eps}}
& \emph{Dialogue - an Android messaging app}\\ 
& \footnotesize{In a team of 7 people we developed an SMS based Android app which allows for encrypted messaging using OpenPGP
and NFC for the key exchange.}\\
\multicolumn{2}{c}{} \\

%------------------------------------------------

\end{tabularx}

%----------------------------------------------------------------------------------------
%	TECHNICAL SKILLS & LANGUAGES
%----------------------------------------------------------------------------------------

\section{Technical Skills \& Languages}
\begin{center}

\vspace{-0.3cm}
\begin{tabularx}{\textwidth}{rX|rl}

& & & \\
\textsc{\large{Programming}} & Scala, Java, Python, C++, Spark, JavaScript, Perl regex &
\textsc{\large{Fluent}} & French, English \\
\textsc{\large{Skills}} & Functional Programming, Compilers, Software Design &
\textsc{\large{Basic}} & German \\

\end{tabularx}

\end{center}

%----------------------------------------------------------------------------------------
%	EXTRA CURRICULAR ACTIVITIES
%----------------------------------------------------------------------------------------

\section{Extra-Curricular Activities}
\begin{tabular}{rl}
\textsc{\large{Undergraduate}} &
Parallel Programming --
Computer Language Processing -- Functional Programming \\
\textsc{\large{Teaching Assistant}} &
Computer Architecture -- Concurrency -- Introduction to Programming (C, C++, Java) \\
\multicolumn{2}{c}{} \\

\textsc{\large{LauzHack Organizer}} &
In a small team, we organised the MLH member hackathon: \href{http://lauzhack.com}{http://lauzhack.com}\\
& Which took place at EPFL with 200 participants and over 50 projects presented\\
\multicolumn{2}{c}{} \\

\end{tabular}

%----------------------------------------------------------------------------------------
%	PERSONAL DATA
%----------------------------------------------------------------------------------------

\section{Personal Data}
\begin{center}
\begin{tabular}{rl}
\textsc{Date of Birth:} & 12 February 1992 \\
\textsc{Nationality:} & Swiss \\
\textsc{Phone:} & +41 77 420 71 06\\
\end{tabular}
\end{center}


\end{document}